\chapter{Introduction}

\section{General Introduction}Agriculture will remain one of the primary sources of livelihood and economic growth in many developing regions. However, farmers will continue to face multiple challenges related to crop health,soil quality, limited access to agricultural information, and difficulty in obtaining farming inputs at the right time. In many rural areas, farming decisions will still be based on experience, manual inspections,or informal advisory suggestions, which will often lead to delays, inaccurate decisions, crop damage,and reduced productivity.

\paragraph{}One of the major challenges in agriculture will be the early identification of plant diseases. Farmers will often fail to detect crop infections in the initial stages, which will allow diseases to spread across fields and cause severe yield loss. Similarly, improper soil nutrient management, lack of soil testing awareness, and limited guidance on crop suitability will negatively affect soil health and crop growth.In addition to this, farmers will have restricted access to agricultural resources such as seeds, fertilizers,and tools, which will require physical visits to marketplaces or intermediaries.
\cite{1}

\paragraph{} Advancements in digital technologies, artificial intelligence, and web-based systems will create opportunities to develop smart agriculture solutions. Machine learning models will enable automated disease detection using plant leaf images, while rule-based and data-driven approaches will support soil analysis and crop recommendations. Chatbot systems will provide instant agricultural guidance,and e-commerce modules will allow convenient access to farming inputs through online platforms.However, most existing solutions will focus on individual functionalities rather than offering an integrated support system for farmers.\cite{2}

\paragraph{}To overcome these gaps, the proposed project “Krushi” will be developed as a unified web-based platform that will integrate plant disease detection, soil analysis, chatbot-based advisory assistance,and an online agricultural shop within a single application. The system will aim to support farmers by improving awareness, assisting decision-making, and promoting smart farming practices. The platform will be designed with a simple, user-friendly interface to ensure accessibility for farmers with varying technical backgrounds.\cite{3}
\paragraph{}The Krushi system will therefore play a significant role in bridging the gap between farmers and modern technology by providing digital agricultural services in an efficient and interactive manner.The platform will eventually contribute toward sustainable agriculture, improved productivity, and technology-driven rural development.\cite{6}




\section{Background}
\paragraph{}Traditional agricultural practices have largely relied on manual observation, experience-based decision-making, and external consultation. Farmers often depend on agricultural experts or local vendors for diagnosing plant diseases and recommending fertilizers. However, this approach is often time-consuming, costly, and not always accessible, especially in rural and remote areas. As a result, delays in identifying crop issues and improper soil management frequently lead to reduced yield and financial instability.\cite{3}\cite{4} 
\paragraph{}In recent years, the adoption of digital technologies in agriculture has increased, giving rise to smart farming solutions that utilize data analysis and intelligent systems. Research in Precision Agriculture has demonstrated that integrating machine learning models, soil analysis techniques, and digital advisory systems can significantly improve farming outcomes. However, most existing solutions focus on individual aspects rather than providing a unified platform. This gap highlights the need for a comprehensive system like Krushi, which will integrate multiple agricultural services into a single, accessible platform for farmers.

\paragraph{}We got the idea of developing a smart agriculture-based web application. That’s why we will develop the “Krushi”. This will be a web-based application which will help farmers to get information and services related to plant disease detection, soil analysis, agricultural guidance, and access to farming products. The system will aim to provide an integrated platform for improving farming efficiency and decision-making.\